332

به حمام مستری وارد شدیم. کاتیا روی زمین زانو زده بود و دستانش را دور گردنش حلقه کرده بود، انگار که داشت خفه می‌شد. برادرش بر روی او خم شده بود و دستگاه تنفسی به دهانش گذاشته بود. مستری چند قدم دورتر ایستاده بود و با نگاهی سرزنش‌آمیز به کاتیا نگاه می‌کرد.

گفتم: «باید آمبولانس خبر کنم؟»

مستری با تمسخر گفت: «او را به خاطر مواد مخدری که در بدنش است دستگیر می‌کنند.»

صبح روز بعد، کورتنی در ساعتی غیرمعمول زود از اتاقش بیرون پرید. لباس خواب «ایجنت پرووکاتوئر» به تن داشت.

اگر او به اندازه‌ای حواس خود را نگه داشته که به مستری چشم‌غره می‌رفت، پس به وضوح در حال مرگ نبود.

پرسید: «چی؟ چه اتفاقی افتاده؟» و خواب را از چشمانش می‌زدود.

سرانجام، وقتی کاتیا از اتاق مستری بیرون آمد، صورتش قرمز و خیس بود. کورتنی دستش را گرفت و او را به سوی مبل در اتاق نشیمن هدایت کرد. کنارش نشست و همچنان دستش را محکم گرفته بود و از تجربیاتش درباره سقط جنین و زیبایی زایمان برایش گفت. به آن زوج نامحتملی که آنجا نشسته بودند نگاه کردم. کورتنی هم فرزند هالیوود بود و هم مادرش.

او همچنین احتمالاً عاقل‌ترین فرد خانه بود. و این ترسناک بود.

راننده‌ای کمی بعد به خانه رسید.

پرسید: «کورتنی کجاست؟»

گفتم: «خواب است.»

گفتند: «او یک ساعت دیگر جلسه دادگاه دارد.»

بر در او زد و وارد شد. اندکی بعد، دسته‌ای از لباس‌ها از اتاق کورتنی به بیرون ریخت، به دنبال آن خود او.

گفت: «باید چیزی برای پوشیدن در دادگاه پیدا کنم» در حالی که انواع لباس‌ها را امتحان می‌کرد و برای دیدن ظاهرشان به آینه حمام می‌رفت و برمی‌گشت. در نهایت، با یک لباس کوتاه کاتیا، عینک هشت دلاری هربال و کتاب «۴۸ قانون قدرت» نوشته رابرت گرین زیر بغل راستش خانه را ترک کرد.

همان روز به خبرنگاران دادگاه گفت: «این یک لباس احمقانه است چون پرونده‌اش احمقانه است.»

در غیاب او، خسارات را بررسی کردیم. در پارچه تخت هربال اثر سوختگی سیگار بود و دیوار پشت در به خاطر کوبیده شدن مداوم خراب شده بود. روی زمین روغن‌های ناشناخته‌ای مالیده بود، شمع‌ها همچنان می‌سوختند و لباس‌هایی همه جا بر روی وسایل نورانی پرتاب شده بود.

در آشپزخانه، تمام درهای یخچال و کمد باز مانده بود. دو شیشه کره بادام زمینی و یک شیشه ژله روی میز بود و درب‌هایشان روی زمین پخش شده بود. گلوله‌های کره بادام زمینی از روی میز، کمدها و قفسه‌های یخچال چکیده بود. به جای باز کردن کیسه‌های نان با استفاده از بستشان، قسمت بالای کیسه‌های پلاستیکی را مثل یک حیوان پاره کرده بود. اصلاً اهمیتی نمی‌داد. گرسنه بود؛ می‌خورد. این هم یکی از ویژگی‌هایی بود که هنرمندان معاشقه تحسین می‌کردند: او می‌توانست به حالت اولیه برگردد. 

وقتی کورتنی از دادگاه برگشت، در کنار گروه خانه نشست و